\documentclass[11pt]{article}
\usepackage[margin=1in]{geometry}
\usepackage{amsmath,amssymb,amsthm,mathtools,mathrsfs}
\usepackage[hidelinks]{hyperref}
\newtheorem{theorem}{Theorem}
\newtheorem{lemma}[theorem]{Lemma}
\newtheorem{proposition}[theorem]{Proposition}
\newtheorem{corollary}[theorem]{Corollary}
\newcommand{\SU}{\operatorname{SU}}
\newcommand{\tr}{\operatorname{tr}}
\newcommand{\Rea}{\operatorname{Re}}
\newcommand{\Var}{\operatorname{Var}}
\newcommand{\dd}{\,\mathrm d}
\title{Exact spatial elimination in the true lattice vacuum:\\marginals, temporal memory, and excitation estimates}
\author{}
\date{8 September 2026}
\begin{document}
\maketitle

\section{The finite interacting Hamiltonian and its actual vacuum}

Let $n\geq2$. Fix a finite lattice with physical oriented links
$\mathsf E$, all its retained oriented plaquette words $\mathsf P$, and
positive coefficients $\kappa_e,b_p$. Reverse traversal uses $U_e^{-1}$.
Take precisely
\begin{equation}\label{eq:H}
 H=-\sum_{e\in\mathsf E}\kappa_e\sum_{a=1}^{n^2-1}X_{e,a}^2
       +\sum_{p\in\mathsf P}b_p(n-\Rea\tr U_p),\qquad
 X_{e,a}f=\left.\frac d{dt}\right|_0f(\ldots,e^{tT_a}U_e,\ldots),
\end{equation}
where $T_a$ is anti-Hermitian, traceless, and
$-2\tr(T_aT_b)=\delta_{ab}$. Its Hilbert space is the probability-Haar
space $L^2(\mathcal Q,\lambda)$ on $\mathcal Q=\SU(n)^{\mathsf E}$.
Every potential term in \eqref{eq:H} is retained. The original metric
$c(X,Y)=-\tfrac12\tr(XY)$ has orthonormal basis $2T_a$ and
$\Delta_e^c=4\sum_aX_{e,a}^2$. Hence the same kinetic operator is
$-\sum_e(\kappa_e/4)\Delta_e^c$.

Let $\psi_0$ be the actual positive normalized ground state,
$H\psi_0=E_0\psi_0$ and $\int\psi_0^2d\lambda=1$. Existence,
smoothness, strict positivity, uniqueness, and gauge invariance hold
for this finite compact elliptic Schr\"odinger operator; they are proved
for \eqref{eq:H} in \texttt{wilson\_finite\_lattice.tex}. Set
\[
 \rho=\psi_0^2,\quad d\mu=\rho d\lambda,\quad
 \mathbb H=L^2(\mathcal Q,\mu),\quad \mathcal Uu=\psi_0u.
\]
The map $\mathcal U:\mathbb H\to L^2(\lambda)$ is unitary, and the
exact ground-state transform is
\begin{align}
 L&=\mathcal U^{-1}(H-E_0)\mathcal U
   =-\rho^{-1}\sum_{e,a}\kappa_eX_{e,a}(\rho X_{e,a}),\label{eq:L}\\
 q(u,v)&=\langle L^{1/2}u,L^{1/2}v\rangle_\mu
   =\sum_{e,a}\kappa_e\int\overline{X_{e,a}u}X_{e,a}v\,d\mu.
   \label{eq:fullform}
\end{align}
Indeed the product rule in $H(\psi_0u)-uH\psi_0$ gives the displayed
divergence operator, and Haar integration by parts gives the form.
These steps use the eigenvalue equation for the full potential.

All inner products below are conjugate-linear in the first argument.
The form domain $\mathcal V$ of $q$ is $H^1(\mathcal Q)$: the fields
span the tangent spaces and the positive smooth $\rho$ is bounded
above and below on the compact product. Its form norm is
$\|u\|_\mathcal V^2=q(u,u)+\|u\|_\mu^2$. Smooth functions are a form
core, and the smooth weighted elliptic operator has domain $H^2$ and
smooth operator core. The compact Sobolev embedding gives compact
resolvent. If $q(u,u)=0$, every derivative is zero and connectedness
gives a constant function. Consequently the kernel is exactly
$\mathbb C1$ and the first nonzero eigenvalue $\delta$ is positive.
These are finite-regulator assertions; $\delta$ is not assumed to have
a positive bound uniform in a spatial continuum limit.

\section{The true marginal, conditional projection, and retained form}

Split the physical links into two disjoint nonempty sets
$\mathsf E=\mathsf S\sqcup\mathsf R$. There is no independence
assumption and no restriction on shared plaquettes in the operator
calculation. Write
\[
 \mathcal Q=\mathcal Q_S\times\mathcal Q_R,\quad U=(s,r),\quad
 d\lambda=d\lambda_Sd\lambda_R,
\]
and define
\begin{equation}\label{eq:marginal}
 \rho_R(r)=\int_{\mathcal Q_S}\rho(s,r)d\lambda_S(s),\qquad
 d\mu_R=\rho_Rd\lambda_R,\qquad
 p(s\mid r)=\frac{\rho(s,r)}{\rho_R(r)}.
\end{equation}
Differentiation under a compact integral proves that $\rho_R$ is smooth.
Its positivity follows from $\rho>0$, and
$\int\rho_Rd\lambda_R=1$ by Fubini. The density $p$ is also smooth
and positive and integrates to one for each $r$.

Let $\mathbb K=L^2(\mathcal Q_R,\mu_R)$ and define
\begin{equation}\label{eq:conditional}
 (Jf)(s,r)=f(r),\qquad
 (J^*u)(r)=\int p(s\mid r)u(s,r)d\lambda_S(s),\qquad
 P=JJ^*,\quad Q=I-P.
\end{equation}
Fubini proves $\|Jf\|_\mu=\|f\|_{\mu_R}$ and
$\langle Jf,u\rangle_\mu=\langle f,J^*u\rangle_{\mu_R}$.
Jensen's inequality proves $\|J^*u\|\leq\|u\|$. Therefore $P$ is
the orthogonal projection onto functions of the retained links, and
\[
 \mathbb H=J\mathbb K\oplus\mathbb H_Q,\qquad
 \mathbb H_Q=\ker J^*=Q\mathbb H.
\]
In particular every $h\in\mathbb H_Q$ is orthogonal to $1=J1$.

For a retained index $i=(e,a)$ with $e\in\mathsf R$, write $X_i=X_{e,a}$
and $\kappa_i=\kappa_e$. These fields act on $r$ and leave $s$ fixed.
Define the real conditional score
\begin{equation}\label{eq:score}
 \ell_i=X_i\log\rho,\qquad
 \bar\ell_i=X_i\log\rho_R,\qquad
 \sigma_i=\ell_i-J\bar\ell_i=X_i\log p(s\mid r).
\end{equation}
Compactness makes every $\sigma_i$ and its derivatives bounded.
Differentiation of \eqref{eq:marginal} gives
$J^*\ell_i=\bar\ell_i$, so $J^*\sigma_i=0$.

\begin{lemma}\label{lem:projection-domain}
Conditional expectation satisfies the exact derivative identity
\begin{equation}\label{eq:conditional-derivative}
 X_iJ^*u=J^*(X_i u)+J^*(\sigma_i u)
\end{equation}
for smooth $u$. It extends boundedly from $\mathcal V$ to
$\mathcal V_R:=H^1(\mathcal Q_R)$, with
\begin{equation}\label{eq:projection-bound}
 a(J^*u,J^*u)\leq2q(u,u)+2\eta\|u\|_\mu^2,
 \quad
 \eta=\sum_{i\in R}\kappa_i\|\sigma_i\|_\infty^2,
\end{equation}
where
\begin{equation}\label{eq:reducedform}
 a(f,g)=\sum_{i\in R}\kappa_i\int\overline{X_if}X_ig\,d\mu_R.
\end{equation}
Both $P$ and $Q$ preserve the full form domain and are bounded for its
norm. The lift $J$ identifies $\mathcal V_R$ with
$\mathcal V\cap J\mathbb K$.
\end{lemma}
\begin{proof}
Differentiate $J^*u=\int p u\,d\lambda_S$ and use $X_ip=p\sigma_i$.
Conditional Cauchy--Schwarz and Fubini give
$\|J^*X_iu\|_{\mu_R}\leq\|X_iu\|_\mu$ and
$\|J^*(\sigma_i u)\|_{\mu_R}\leq\|\sigma_i\|_\infty\|u\|_\mu$.
Squaring their sum with $(x+y)^2\leq2x^2+2y^2$ and summing all
retained $\kappa_i$ proves \eqref{eq:projection-bound}.
Smooth approximation extends the derivative identity in the weak
Sobolev sense and the norm bound to $\mathcal V$. Since
$q(Jf,Jf)=a(f,f)$ and $J$ is an isometry in $L^2$, this proves the
claims about $P$ and $Q$. Conversely the weak retained derivatives of
$Jf$ are the lifts of those of $f$, so $Jf\in\mathcal V$ implies
$f\in\mathcal V_R$. Repeated differentiation of the compact conditional
integral also shows $J^*:H^k(\mathcal Q)\to H^k(\mathcal Q_R)$ is
bounded for each nonnegative integer $k$; in particular these maps
preserve smooth functions.
\end{proof}

The form $a$ is closed on $\mathcal V_R$ and has the self-adjoint
weighted elliptic generator
\begin{equation}\label{eq:A}
 A=-\rho_R^{-1}\sum_{i\in R}\kappa_iX_i(\rho_RX_i),
 \qquad D(A)=H^2(\mathcal Q_R).
\end{equation}
For a nonconstant smooth retained observable $F$, let
$f=F-\mu_R(F)$ and $u_F=\psi_0Jf$. Then
\begin{equation}\label{eq:Rayleigh}
 \langle\psi_0,u_F\rangle_\lambda=0,\qquad
 \frac{\langle u_F,(H-E_0)u_F\rangle_\lambda}{\|u_F\|_\lambda^2}
 =\frac{\displaystyle\sum_{i\in R}\kappa_i\mu_R(|X_iF|^2)}
        {\displaystyle\mu_R(|F|^2)-|\mu_R(F)|^2}.
\end{equation}
The numerator follows from \eqref{eq:fullform}, since the eliminated-link
derivatives of $Jf$ vanish; the denominator and orthogonality follow
from Fubini. Positivity of $\rho_R$ makes the denominator strictly
positive. Eliminated kinetic coefficients and all magnetic coefficients
still affect $\psi_0$, $\rho_R$, and the conditional density.

For the original gauge law $U_e\mapsto h_{s(e)}^{-1}U_eh_{t(e)}$,
the same vertex transformation acts separately on $s$ and $r$.
Gauge invariance of $\rho$ and Haar substitution in $s$ prove
$\rho_R(h\cdot r)=\rho_R(r)$ and covariance of $p(s\mid r)$.
Thus $J,J^*,P,Q$ intertwine these gauge actions. The equally valid law
$g_{s(e)}U_eg_{t(e)}^{-1}$ uses $g_v=h_v^{-1}$ and gives exactly the
same subspaces. Every formula therefore restricts to gauge-invariant
observables and vectors. For such a nonconstant $F$, \eqref{eq:Rayleigh}
is an upper bound on the physical finite-lattice gap. The infimum of
these quotients over retained invariant functions is at least the
infimum over all full invariant functions; a lower bound on one
retained quotient alone is not a full spectral lower bound.

\section{The exact coupling to eliminated functions and its domains}

For $f\in C^\infty(\mathcal Q_R)$, expand \eqref{eq:L} and compare
with \eqref{eq:A}. The retained second-derivative terms cancel, and
all eliminated derivatives of $Jf$ are zero. Hence
\begin{equation}\label{eq:B}
 LJf=JAf+Bf,\qquad
 Bf=-\sum_{i\in R}\kappa_i\sigma_i J(X_if),\qquad J^*Bf=0.
\end{equation}
This is the exact first-order map from retained variation to the
eliminated subspace. It extends boundedly
$B:\mathcal V_R\to\mathbb H_Q$, because
\begin{equation}\label{eq:Bbound}
 \|Bf\|_\mu^2\leq\eta\,a(f,f).
\end{equation}
Pointwise Cauchy--Schwarz with the factors $\sqrt{\kappa_i}\sigma_i$
and $\sqrt{\kappa_i}X_if$, followed by integration, proves this bound.
The exact squared norm retains the full conditional covariance:
\begin{equation}\label{eq:Fisher}
 \|Bf\|_\mu^2
 =\int\sum_{i,j\in R}\kappa_i\kappa_j
       I_{ij}(r)\overline{X_if(r)}X_jf(r)\,d\mu_R(r),
 \quad I_{ij}=J^*(\sigma_i\sigma_j).
\end{equation}
The matrix $I(r)$ is positive semidefinite by the exact identity
\[
 \sum_{i,j}\overline z_i I_{ij}(r)z_j
 =J^*\left(\left|\sum_i z_i\sigma_i\right|^2\right)(r)\geq0.
\]
Because $J^*\ell_i=\bar\ell_i$, it is exactly the conditional covariance
of $X_i\log\rho$ and $X_j\log\rho$ under $p(s\mid r)$.

Define $\mathcal V_Q=\mathcal V\cap\mathbb H_Q$ and restrict the
full form:
\begin{equation}\label{eq:Cform}
 c_Q(h,k)=q(h,k),\qquad h,k\in\mathcal V_Q.
\end{equation}
This is a closed densely defined nonnegative form on $\mathbb H_Q$.
Closedness follows since $\mathbb H_Q$ is $L^2$ closed. Density and a
form core follow by applying the bounded $Q$ of Lemma
\ref{lem:projection-domain} to smooth approximations; thus
$Q C^\infty(\mathcal Q)$ is a form core. Let $C$ be its associated
self-adjoint operator. There is no unsupported assertion here that
the formal differential compression $QLQ$ has an automatically chosen
operator domain. The form determines the domain uniquely by
\[
 D(C)=\{h\in\mathcal V_Q:\text{there exists }v\in\mathbb H_Q
       \text{ with }c_Q(k,h)=\langle k,v\rangle
       \text{ for every }k\in\mathcal V_Q\},\quad Ch=v.
\]
For smooth $h\in\mathbb H_Q$ this definition gives $Ch=QLh$ by
integration by parts. Since every $h\in\mathbb H_Q$ is orthogonal to
$1$, the full Poincare inequality gives
\begin{equation}\label{eq:Cgap}
 c_Q(h,h)\geq\delta\|h\|_\mu^2,
 \qquad C\geq\delta I>0.
\end{equation}
In particular $C^{-1}$ exists as a bounded operator. It and all
$(C+z)^{-1}$, $z\geq0$, map into $D(C)\subset\mathcal V_Q$.

For $f,g\in\mathcal V_R$ and $h,k\in\mathcal V_Q$, the exact block
form is
\begin{equation}\label{eq:blockform}
 q(Jf+h,Jg+k)=a(f,g)+\langle Bf,k\rangle
                         +\langle h,Bg\rangle+c_Q(h,k).
\end{equation}
For smooth functions this follows from \eqref{eq:B}; the bounded map
$B:\mathcal V_R\to\mathbb H_Q$ and the form-domain decomposition in
Lemma \ref{lem:projection-domain} extend it to the displayed domains.

The adjoint coupling has an explicit expression as well. For smooth
$h\in\mathbb H_Q$, Haar integration by parts in $r$ gives
\begin{equation}\label{eq:Bstar}
 B^*h=\rho_R^{-1}\sum_{i\in R}\kappa_i
                   X_i\bigl(\rho_R J^*(\sigma_i h)\bigr).
\end{equation}
Indeed $\langle Bf,h\rangle=-\sum_i\kappa_i
\int\rho_R\overline{X_if}J^*(\sigma_ih)d\lambda_R$, which becomes
$\langle f,B^*h\rangle$ with the sign in \eqref{eq:Bstar}.
Using \eqref{eq:conditional-derivative} with $u=\sigma_i h$ expands it
as the finite sum
\[
 B^*h=\sum_{i\in R}\kappa_i
 \left[J^*(\sigma_iX_ih)+J^*((X_i\sigma_i)h)
       +J^*(\sigma_i^2h)+\bar\ell_iJ^*(\sigma_i h)\right].
\]
Every coefficient is bounded, so this formula extends to a bounded map
$B^*:\mathcal V_Q\to\mathbb K$. On $\mathbb H_Q$ without this
regularity, the symbol $B^*$ is instead the bounded dual map into
$\mathcal V_R^*$ specified by $\langle f,B^*h\rangle=\langle Bf,h\rangle$.
These two definitions agree on their common domain.

Replacing $X_i$ by $X_i^c=2X_i$, $\sigma_i$ by $2\sigma_i$, and every
kinetic coefficient by $\kappa_i/4$ leaves \eqref{eq:A}, \eqref{eq:B},
and \eqref{eq:blockform} unchanged. This records the full original
metric dictionary for the eliminated coupling as well as for $H$.

\section{Exact compressed resolvents as closed Schur forms}

For $z>0$, put $R_z=(C+z)^{-1}$ and define on $\mathcal V_R$
\begin{equation}\label{eq:Schurform}
 s_z(f,g)=a(f,g)+z\langle f,g\rangle-\langle Bf,R_zBg\rangle.
\end{equation}
The last term is well defined by \eqref{eq:Bbound}. All occurrences of
$B^*R_zB$ below mean this form on the stated domain, not a product of
unbounded operators with unspecified domains.

\begin{theorem}\label{thm:Schur}
The form $s_z$ is closed, positive, and coercive on $\mathcal V_R$.
Its self-adjoint operator $S_z$ obeys $S_z\geq zI$ and
\begin{equation}\label{eq:Schurresolvent}
 J^*(L+z)^{-1}J=S_z^{-1}
   =\left[A+z-B^*(C+z)^{-1}B\right]^{-1}.
\end{equation}
The eliminated component of $(L+z)^{-1}Jg$ is exactly
$-R_zB S_z^{-1}g$. Moreover, in operator order,
\begin{equation}\label{eq:resolventorder}
 J^*(L+z)^{-1}J\ \geq\ (A+z)^{-1}.
\end{equation}
\end{theorem}
\begin{proof}
The block form and completion of the square give, for every
$f\in\mathcal V_R$, $h\in\mathcal V_Q$,
\begin{equation}\label{eq:square}
 q(Jf+h,Jf+h)+z(\|f\|^2+\|h\|^2)
 =s_z(f,f)+\|(C+z)^{1/2}(h+R_zBf)\|^2.
\end{equation}
Thus $s_z(f,f)$ is the infimum over $h$, uniquely attained at
$h=-R_zBf$, and nonnegativity of $q$ gives $s_z(f,f)\geq z\|f\|^2$.
For an explicit form-norm lower bound, apply
\eqref{eq:projection-bound} to $v=Jf+h$, for which $J^*v=f$:
\[
 a(f,f)+\|f\|^2
 \leq2q(v,v)+(2\eta+1)\|v\|^2
 \leq K_z\bigl(q(v,v)+z\|v\|^2\bigr),
 \quad K_z=\max\{2,(2\eta+1)/z\}.
\]
Take the minimizing $h$ to obtain
$s_z(f,f)\geq K_z^{-1}(a(f,f)+\|f\|^2)$.
Conversely $s_z(f,f)\leq a(f,f)+z\|f\|^2$ by choosing $h=0$.
These two bounds make its norm equivalent to the complete
$\mathcal V_R$ norm and prove closedness and coercivity. The
associated operator therefore exists and is at least $zI$.

For $g\in\mathbb K$, let $v=(L+z)^{-1}Jg$ and decompose it as
$v=Jf+h$ using the bounded form projection. The weak resolvent
equation tested against $k\in\mathcal V_Q$ gives
$c_Q(k,h)+z\langle k,h\rangle=-\langle k,Bf\rangle$.
The definition of $C$ proves $h=-R_zBf$. Testing against $Jf'$ for
$f'\in\mathcal V_R$ then gives
$s_z(f',f)=\langle f',g\rangle$, so $f=S_z^{-1}g$.
This proves both the compressed inverse and the reconstruction formula
with their domains resolved.

Finally, $s_z(f,f)\leq a(f,f)+z\|f\|^2$. For a positive coercive
closed form $t$ with associated operator $T$, completing the square in
its form space gives
$\langle g,T^{-1}g\rangle=\sup_f[2\Rea\langle g,f\rangle-t(f,f)]$.
Applying this formula to the two forms proves
\eqref{eq:resolventorder}. It is an order statement for the resolvents;
no ordering of their time exponentials is inferred from it.
\end{proof}

This is a weighted, infinite-rank Schur complement. The projection and
effective-resolvent method itself is established in Feshbach's work and
in the modern Feshbach--Schur formulation \cite{Feshbach,DSS}. The proof
above supplies the conditional-score coupling, form domains, and
coercivity needed for this particular infinite-dimensional projection.

\section{Full temporal correlations and exact memory}

Define the compressed correlation operator on the retained Hilbert space
by
\begin{equation}\label{eq:Kt}
 K(t)=J^*e^{-tL}J,\qquad t\geq0.
\end{equation}
It is a positive self-adjoint contraction with $K(0)=I$ and $K(t)1=1$.
These statements follow respectively from positivity in the spectral
sense, contraction, the isometry $J$, and $L1=0$. Its matrix elements
are the actual full quantum vacuum correlations, since
\begin{equation}\label{eq:quantumcorrelation}
 \langle f,K(t)g\rangle_{\mu_R}
 =\langle\psi_0Jf,e^{-t(H-E_0)}\psi_0Jg\rangle_\lambda.
\end{equation}
The spectral theorem also gives its norm-convergent Laplace integral
\begin{equation}\label{eq:Laplace}
 \int_0^\infty e^{-zt}K(t)\,dt
 =J^*(L+z)^{-1}J=S_z^{-1},\qquad z>0.
\end{equation}
The integral is convergent in operator norm because $\|K(t)\|\leq1$.

\begin{theorem}\label{thm:memory}
For smooth initial $f_0\in C^\infty(\mathcal Q_R)$, set
$f(t)=K(t)f_0$. It satisfies the exact equation in $\mathbb K$
\begin{equation}\label{eq:memory}
 f'(t)=-Af(t)+\int_0^t B^*e^{-(t-s)C}Bf(s)\,ds,
 \qquad f(0)=f_0.
\end{equation}
The integral is a well-defined Bochner integral; its positive sign is
fixed by eliminating the other component of the full equation.
For every smooth $f$,
\begin{equation}\label{eq:secondmoment}
 \langle f,K(t)f\rangle
 =\|f\|^2-t\,a(f,f)
       +\frac{t^2}{2}\bigl(\|Af\|^2+\|Bf\|^2\bigr)+o(t^2).
\end{equation}
The equality $K(t)=e^{-tA}$ for all $t\geq0$ holds if and only if
$B=0$ on $\mathcal V_R$.
\end{theorem}
\begin{proof}
Let $v(t)=e^{-tL}Jf_0$ and $h(t)=Qv(t)$. Smooth initial data belong to
all powers of the compact elliptic operator $L$, so spectral calculus
and elliptic regularity make $v(t)$ smooth in space and continuously
differentiable in every required Sobolev norm for $t\geq0$. Lemma
\ref{lem:projection-domain} implies the same spatial regularity for
$f(t)=J^*v(t)$ and $h(t)$. The full equation $v'=-Lv$ and
\eqref{eq:blockform} therefore give
\[
 f'=-Af-B^*h,\qquad h'=-Bf-Ch,\qquad h(0)=0.
\]
For the second equation, differentiating
$e^{-(t-s)C}h(s)$ and integrating in $s$ gives
\begin{equation}\label{eq:eliminatedevolution}
 h(t)=-\int_0^te^{-(t-s)C}Bf(s)\,ds.
\end{equation}
This differentiation is justified since the smooth $h(s)$ lies in
$D(C)$ with $Ch(s)=QLh(s)$ continuously in $\mathbb H_Q$.

The map $B^*:\mathcal V_Q\to\mathbb K$ was proved bounded above.
Spectral calculus for $C\geq0$ gives
\[
 \|e^{-\tau C}\|_{\mathbb H_Q\to\mathcal V_Q}
 \leq1+(2e\tau)^{-1/2},\qquad\tau>0,
\]
because $\sup_{x\geq0}\sqrt{x}e^{-\tau x}=(2e\tau)^{-1/2}$.
Consequently the integrand in \eqref{eq:eliminatedevolution} is
integrable in the form norm near $s=t$, as $Bf(s)$ is continuous in
$\mathbb H_Q$. Applying the bounded $B^*$ to that form-norm integral
and substituting into the first equation proves \eqref{eq:memory}
as a strong $\mathbb K$ identity. This also specifies rigorously the
memory-kernel domain: $B^*e^{-\tau C}B$ maps $\mathcal V_R$ to
$\mathbb K$ for $\tau>0$, with an integrable $O(1+\tau^{-1/2})$
bound on bounded time intervals. This is an exact projection-memory
construction of the type introduced by Mori \cite{Mori}, with the
Euclidean generator and signs derived here.

For \eqref{eq:secondmoment}, apply the scalar spectral theorem to the
smooth vector $Jf$, which lies in $D(L^2)$. The zeroth, first and
second moments are $\|f\|^2$, $a(f,f)$, and
\[
 \|LJf\|_\mu^2=\|JAf+Bf\|_\mu^2=\|Af\|_{\mu_R}^2+\|Bf\|_\mu^2.
\]
The second-order Taylor remainder is $o(t^2)$ by dominated convergence
against its finite second spectral moment. This proves the expansion.
The corresponding expansion for $e^{-tA}$ has second coefficient
$\|Af\|^2/2$. Equality of all the correlation operators therefore
forces $Bf=0$ for every smooth $f$, and density and \eqref{eq:Bbound}
extend this to $\mathcal V_R$. Conversely, if $B=0$, the exact closed
form \eqref{eq:blockform} is the orthogonal direct sum of $a$ and $c_Q$.
Its self-adjoint operator and semigroup are direct sums as well, giving
$K(t)=e^{-tA}$.
\end{proof}

In particular the explicit second-derivative diagnostic is
\[
 \left.\frac{d^2}{dt^2}\right|_{t=0}
 \left[\langle Jf,e^{-tL}Jf\rangle_\mu
              -\langle f,e^{-tA}f\rangle_{\mu_R}\right]=\|Bf\|_\mu^2.
\]

For the unrestricted retained function space this condition has a
concrete statistical characterization:
\begin{equation}\label{eq:factorization}
 B=0\text{ on }\mathcal V_R
 \quad\Longleftrightarrow\quad
 \rho(s,r)=\rho_S(s)\rho_R(r)
 \text{ for positive probability densities }\rho_S,\rho_R.
\end{equation}
The reverse implication follows immediately from $\sigma_i=0$.
For the forward implication, \eqref{eq:B} vanishes for every smooth
test function $f$. The resulting functions are smooth, so vanishing
in $L^2$ means vanishing pointwise. At any fixed $r$, a smooth function
with arbitrary prescribed differential at $r$ is obtained from local
coordinates and a smooth cutoff. The independent tangent basis $X_i$
and positive $\kappa_i$ then force each $\sigma_i(s,r)=0$. Thus every
$r$ derivative of $p(s\mid r)$ is zero, and connectedness of
$\mathcal Q_R$ makes it independent of $r$. Normalization gives the
factorization. If one restricts tests to gauge-invariant functions,
the exact criterion is $B=0$ on that invariant form space; the stronger
factorization is not asserted from that smaller test space.

\section{Constructed excitation vectors using the eliminated resolvent}

The marginal quotient can be improved by an explicit full-state vector.
Let $0\ne f\in\mathcal V_R$ satisfy $\mu_R(f)=0$. For every $z\geq0$
define
\begin{equation}\label{eq:correctedstate}
 w_z=(C+z)^{-1}Bf,\qquad
 u_z=\psi_0(Jf-w_z).
\end{equation}
Equation \eqref{eq:Cgap} guarantees existence even at $z=0$, and
$w_z\in\mathcal V_Q$. Orthogonality gives
$\langle\psi_0,u_z\rangle_\lambda=0$ and
$\|u_z\|_\lambda^2=\|f\|_{\mu_R}^2+\|w_z\|_\mu^2>0$.
All gauge actions commute with the forms and their resolvents; hence
if $f$ is invariant, $u_z$ is invariant as well.
For a full original-Hilbert-space form-domain vector $u$, write
\[
 \mathfrak h_0[u]=\|(H-E_0)^{1/2}u\|_\lambda^2
                =q(u/\psi_0,u/\psi_0).
\]
This is the energy used below. It equals
$\langle u,(H-E_0)u\rangle$ when $u\in D(H)$, but no such stronger
domain membership is assumed for the constructed form-domain vector.

\begin{theorem}\label{thm:improvedtrial}
The actual full-Hamiltonian Rayleigh quotient of \eqref{eq:correctedstate}
is
\begin{equation}\label{eq:improvedquotient}
 \frac{\mathfrak h_0[u_z]}{\|u_z\|^2}
 =\frac{a(f,f)-\langle Bf,w_z\rangle-z\|w_z\|^2}
        {\|f\|^2+\|w_z\|^2}
 \leq\frac{a(f,f)}{\|f\|^2}.
\end{equation}
The inequality is strict whenever $Bf\ne0$. The numerator is
nonnegative. For $z=0$, $-w_0$ uniquely minimizes the full Dirichlet
energy among all $Jf+h$, $h\in\mathcal V_Q$, with the retained
conditional mean $f$ fixed.
\end{theorem}
\begin{proof}
The equation $(C+z)w_z=Bf$ gives
$c_Q(w_z,w_z)=\langle Bf,w_z\rangle-z\|w_z\|^2$.
Insert $h=-w_z$ in \eqref{eq:blockform}. The two cross terms subtract
$2\langle Bf,w_z\rangle$, yielding exactly the numerator in
\eqref{eq:improvedquotient}. It is nonnegative because it is the full
closed form $q(Jf-w_z,Jf-w_z)$. The scalar
$\langle Bf,w_z\rangle=\langle Bf,(C+z)^{-1}Bf\rangle$ is
nonnegative and strictly positive when $Bf\ne0$. The numerator is
therefore at most $a(f,f)$ and the denominator at least $\|f\|^2$;
when $Bf\ne0$ the numerator is strictly smaller. The base energy is
positive for a nonzero centered $f$, because $\rho_R>0$ on a connected
space. This proves the claimed strict inequality. Completion of the
square in \eqref{eq:square} with $z=0$ remains valid because $C>0$,
and proves the minimization statement.
\end{proof}

For a centered retained $f$, define the actual excitation spectral
probability measure
\[
 \nu_f(E)=\|f\|^{-2}
 \langle\psi_0Jf,\mathbf1_E(H-E_0)\psi_0Jf\rangle.
\]
Its first moment is $a(f,f)/\|f\|^2$, its second moment for smooth
$f$ is $(\|Af\|^2+\|Bf\|^2)/\|f\|^2$, and its Laplace transform
is $\langle f,K(t)f\rangle/\|f\|^2$. The corrected states have their
own actual spectral measures with first moment
\eqref{eq:improvedquotient}. In particular each invariant corrected
state gives a genuine upper bound on the full physical spectral gap;
none of these energies is assigned from an independently chosen
single-slice density.

\section{Exact nonlinear Wilson integration leading to this marginal}

We now construct the connection to the completed compact-link
integration. Specialize to $\SU(2)$ and the fixed spatial regulator
with the original action metric, spatial mesh $a>0$, and coupling
$g_{\rm YM}>0$. Retain
\begin{equation}\label{eq:physicalcoefficients}
 \kappa_e=\frac{2g_{\rm YM}^2}{a},\qquad
 b_p^{\rm Hamiltonian}=\frac1{2g_{\rm YM}^2a},\qquad
 b_t(\epsilon)=\frac{a}{2g_{\rm YM}^2\epsilon},\qquad
 b_s(\epsilon)=\frac{\epsilon}{2g_{\rm YM}^2a}.
\end{equation}
The companion \texttt{temporal\_transfer.tex} defines the positive
normalized transfer operator
\[
 F_\epsilon=e^{-\epsilon V/2}K_\epsilon e^{-\epsilon V/2}.
\]
The exact temporal link kernel has
spin-$j$ multiplier $I_{2j+1}(2b_t)/I_1(2b_t)$, temporal vertex
integration gives the Gauss projector $P_G$, and the raw transfer is
\begin{equation}\label{eq:rawrelation}
 T_\epsilon^{\rm raw}=c_\epsilon F_\epsilon P_G,
 \qquad c_\epsilon=[e^{-2b_t}Z(b_t)]^{|\mathsf E|},\qquad
 Z(b)=I_1(2b)/b.
\end{equation}
That companion proves, on the actual dense invariant core and then
strongly on the full Hilbert space,
$F_{T/m}^{\,m}\to e^{-TH}$ for fixed $T>0$. This is a time limit at
fixed spatial graph, with every coefficient in
\eqref{eq:physicalcoefficients} retained.

For $T>0$ and integer $m\geq1$, put $\epsilon=T/m$ and
\begin{equation}\label{eq:slabdensities}
 \phi_{T,m}=F_{T/m}^{\,m}1,\qquad
 \rho_{T,m}(U)=\frac{\phi_{T,m}(U)^2}{\|\phi_{T,m}\|_\lambda^2},\qquad
 \rho_{R;T,m}(r)=\int\rho_{T,m}(s,r)d\lambda_S(s).
\end{equation}
All these functions are positive, and $1$ is gauge invariant, so the
Gauss projector is retained automatically. Kernel composition and
Fubini identify \eqref{eq:slabdensities} with the central-slice density
of a two-sided finite Euclidean Wilson slab with slices
$-m,-m+1,\ldots,m$, fixed central slice $U$, and constant outer
boundary states. Each temporal plaquette has weight $b_t$, every
interior spatial plaquette has weight $b_s$, and spatial plaquettes
on the two outer boundary slices have half weight. The central
magnetic factor is full, because the two half factors multiply.
The two raw factors $c_\epsilon^m$ multiply to $c_\epsilon^{2m}$ in
both numerator and normalization and cancel exactly. This proves a
specific Euclidean-integral representation, including both temporal
sides and their boundary weights.

Assume now that $\mathsf S$ is the selected spatial-link set with no
two selected edges belonging to one spatial plaquette, and that each
plaquette uses any one physical edge at most once. Apply the selection
only at the central time slice. A temporal plaquette has only one
spatial edge at that slice, so no spacetime plaquette contains two
selected central links either. Let $Y$ denote every remaining slab
variable except the retained central links $r$.
For each selected central edge $e$, cyclically rotate and, when
necessary, invert a plaquette inside its real trace to write all its
incident contributions as $\Rea\tr(U_eV_{ep})$. Define
\begin{align}
 H_e(Y,r)&=\sum_{p\ni(e,0)}b_pV_{ep},&
 B_e&=\sum_{p\ni(e,0)}b_p,\nonumber\\
 r_e^2(Y,r)&=\tfrac12\tr(H_eH_e^\dagger)
 =\sum_{p\ni(e,0)}b_p^2
   +\sum_{p<q\ni(e,0)}b_pb_q\Rea\tr(V_{ep}V_{eq}^\dagger).
   \label{eq:allstaples}
\end{align}
The weights in this sum are the actual anisotropic weights of their
plaquettes. For a spatial interior edge there are four spatial staples
of weight $b_s$ and two temporal staples of weight $b_t$, so
$B_e=4b_s+2b_t$. At a spatial boundary the displayed incident sum,
rather than this interior count, is used.

The exact single-link integral is
\begin{equation}\label{eq:Zintegration}
 \int_{\SU(2)}e^{-2B_e+\Rea\tr(U_eH_e)}dU_e
   =e^{-2B_e}Z(r_e),\qquad
 Z(r)=\sum_{k=0}^\infty\frac{r^{2k}}{k!(k+1)!},\quad Z(0)=1.
\end{equation}
To prove it, write every $\SU(2)$ matrix as a unit real quaternion.
The real linear combination $H_e$ then obeys
$H_eH_e^\dagger=r_e^2I$ and $\det H_e=r_e^2$.
For $r_e>0$, $H_e/r_e\in\SU(2)$ and Haar translation reduces the
integral to
$(2/\pi)\int_{-1}^1e^{2r_ex}\sqrt{1-x^2}dx$.
Termwise integration of its uniformly convergent exponential series,
using
$\int_{-1}^1x^{2k}\sqrt{1-x^2}dx
=\pi(2k)!/[2\cdot4^kk!(k+1)!]$, gives
\eqref{eq:Zintegration}. The even moment formula follows by integration
by parts from the semicircle area at $k=0$. At $r_e=0$ the quaternion
is zero and the integral is exactly one before its constant factor.

No incident plaquette is shared by selected links. Thus the exact
central retained density in \eqref{eq:slabdensities} is
\begin{equation}\label{eq:integratedslab}
 \rho_{R;T,m}(r)=
 \frac{\displaystyle\int
    e^{-S_{\rm rest}(Y,r)}
    \prod_{e\in\mathsf S}e^{-2B_e}Z(r_e(Y,r))\,dY}
      {\displaystyle\int_{\mathcal Q_R}\int
    e^{-S_{\rm rest}(Y,r')}
    \prod_{e\in\mathsf S}e^{-2B_e}Z(r_e(Y,r'))\,dY\,d\lambda_R(r')}.
\end{equation}
Here $S_{\rm rest}$ is precisely the slab action with those incident
plaquettes removed; it includes every other plaquette and the stated
outer boundary weights. The integrals over $Y$ include all temporal
vertex links. Formula \eqref{eq:integratedslab} follows by applying
\eqref{eq:Zintegration} to the finite positive integrand and Fubini.
All cross-staple traces, the two temporal staples, and all induced
interactions remain in it. For a general split $\mathsf S$, the
unfactorized exact integral over selected links in
\eqref{eq:slabdensities} remains valid; only the product expression
\eqref{eq:integratedslab} uses the no-shared-plaquette property.

\begin{theorem}\label{thm:slablimit}
The exact true-vacuum marginal is the sequential $L^1$ limit
\begin{equation}\label{eq:sequentiallimit}
 \rho_R=\lim_{T\to\infty}\lim_{m\to\infty}\rho_{R;T,m},
 \qquad\epsilon=T/m,
\end{equation}
with the spatial graph and $a,g_{\rm YM}$ fixed. In particular the
nonlinear integration \eqref{eq:integratedslab}, performed with all
its interactions, has this precise relation to the true vacuum.
\end{theorem}
\begin{proof}
Strong transfer convergence gives
$\phi_{T,m}\to\phi_T:=e^{-TH}1$ in $L^2(\lambda)$ for each fixed
$T>0$. This nonzero positive vector has a nonzero norm. Normalization
is continuous at a nonzero vector, and
$\|v^2-w^2\|_1\leq(\|v\|_2+\|w\|_2)\|v-w\|_2$ for real
functions. Consequently
\[
 \rho_{T,m}\longrightarrow
 \rho_T:=\frac{(e^{-TH}1)^2}{\|e^{-TH}1\|^2}
 \quad\text{in }L^1(\lambda).
\]
Marginalization is an $L^1$ contraction by Fubini and the triangle
inequality, so the retained marginals converge in $L^1(\lambda_R)$.

Set $c_0=\langle\psi_0,1\rangle_\lambda>0$. Spectral decomposition,
the unique ground state, and the finite gap $\delta$ give
\[
 e^{TE_0}e^{-TH}1=c_0\psi_0+r_T,\qquad
 r_T\perp\psi_0,\qquad
 \|r_T\|\leq e^{-T\delta}\sqrt{1-c_0^2}.
\]
The norm of this vector is $\sqrt{c_0^2+\|r_T\|^2}$. Subtracting
$\psi_0$ after normalizing and using the triangle inequality gives
\[
 \left\|\frac{e^{-TH}1}{\|e^{-TH}1\|}-\psi_0\right\|_2
 \leq\frac{2\|r_T\|}{c_0}.
\]
Thus the same square-density inequality and marginal contraction yield
\begin{equation}\label{eq:vacuumerror}
 \|\rho_{R;T}-\rho_R\|_1
 \leq\|\rho_T-\rho\|_1
 \leq\frac{4\sqrt{1-c_0^2}}{c_0}e^{-T\delta}.
\end{equation}
This proves \eqref{eq:sequentiallimit}. It does not exchange the two
limits or infer a spatial continuum limit from them.
\end{proof}

The resulting excitation estimate can be obtained continuously from
these exact marginal integrations. Fix a nonconstant smooth invariant
$F(r)$ and put
\[
 \gamma_R(F)=\sum_{i\in R}\kappa_i|X_iF|^2,
 \qquad Q_R(F)=\frac{\mu_R(\gamma_R(F))}{\Var_{\mu_R}(F)}.
\]
For a probability density $\widetilde\rho_R$ let $\widetilde Q_R(F)$
be the same ratio with that density. Set
$\eta_1=\|\widetilde\rho_R-\rho_R\|_1$,
$M_F=\|F\|_\infty$, $G_F=\|\gamma_R(F)\|_\infty$, and
$v_F=\Var_{\mu_R}(F)>0$. Direct integration gives
\[
 |\widetilde\mu_R(\gamma_R)-\mu_R(\gamma_R)|\leq G_F\eta_1,
 \quad
 |\Var_{\widetilde\mu_R}(F)-v_F|\leq3M_F^2\eta_1.
\]
The second inequality uses the second-moment error at most
$M_F^2\eta_1$ and
$\bigl|\lvert\widetilde\mu_R(F)\rvert^2-\lvert\mu_R(F)\rvert^2\bigr|
\leq2M_F^2\eta_1$ for the mean term. If
$\eta_1\leq v_F/(6M_F^2)$, the new variance is at least $v_F/2$;
subtracting the two fractions then proves the explicit estimate
\begin{equation}\label{eq:quotienterror}
 |\widetilde Q_R(F)-Q_R(F)|
 \leq\frac{2G_F}{v_F}\eta_1
      +\frac{6G_FM_F^2}{v_F^2}\eta_1.
\end{equation}
In particular the sequential limits of the quotients constructed from
\eqref{eq:integratedslab} equal the genuine full-vacuum Rayleigh
quotient \eqref{eq:Rayleigh}, and hence its actual spectral first
moment. The controlled finite integrals, the true marginal, and the
quantum excitation estimate are related by the displayed maps and
limits. The single-time density proportional to
$\exp[-\sum_pb_p(2-\Rea\tr U_p)]$ is nowhere substituted for
$\psi_0^2$.

\section{Verification and literature provenance}

The calculation uses the full finite interacting Hamiltonian, not a
Gaussian or Abelian replacement. The map $B$, its conditional
covariance, the closed Schur form, the memory integral, and the
corrected full trial states all retain the actual $\rho=\psi_0^2$.
The scalar Schur/memory signs and corrected Rayleigh algebra are also
checked in \texttt{check\_spatial\_elimination.py}; that finite exact
algebra check does not replace the proved form-domain argument.
It passed 12 exact scalar/block assertions, together with exact
matrix-entry checks. An additional explicit rational-matrix quotient
checks the strict decrease; it is an algebra diagnostic and is not
presented as a lattice vacuum or a continuum approximation.

The general projection and Schur ideas are established methods, not
claimed as new results here. The locally retained source of
Dusson--Sigal--Stamm was inspected at Theorem 1.2 and the subsequent
effective-interaction equations; the retained projection here has
infinite rank, so its bounded finite-rank formulas are not invoked
without the form proof above. The earlier complete finite-vacuum and
time-transfer sources in this directory provide the retained kinetic,
potential, metric, and temporal-normalization dictionaries used here.
No infinite-volume or spatial-continuum spectral conclusion is made.

\begin{thebibliography}{9}
\bibitem{Feshbach}
H. Feshbach, \emph{Unified Theory of Nuclear Reactions},
Reviews of Modern Physics \textbf{36} (1964), 1076--1078,
\href{https://doi.org/10.1103/RevModPhys.36.1076}
{doi:10.1103/RevModPhys.36.1076}.

\bibitem{DSS}
G. Dusson, I. M. Sigal, and B. Stamm,
\emph{The Feshbach--Schur map and perturbation theory},
\href{https://arxiv.org/abs/2105.02058}{arXiv:2105.02058v1} (2021),
Theorem 1.2 and equations (1.10)--(1.16).

\bibitem{Mori}
H. Mori, \emph{Transport, Collective Motion, and Brownian Motion},
Progress of Theoretical Physics \textbf{33} (1965), 423--455,
\href{https://doi.org/10.1143/PTP.33.423}{doi:10.1143/PTP.33.423}.
\end{thebibliography}
\end{document}
